
% Why are organic (and inorganic) solar cells important.
In x year, solar made up x\% of energy produced in the world and y\% energy produced in Texas.
Solar is an increasingly valuable and important energy source, and a highreturn green energy source.
Solar prduction has increased x\% between 2010 and 2026.
Solar often has a return on cost within x\% years in Texas.

Organic Solar Cell (OSC) research has experienced incredible advancements over the previous 12 years \autocite{xuChallengingPM6likeDonor2023}.
Power Conversion Efficiency (PCE) of modern OSCs produced in labs have exceeded 20\%, comparable to traditional silicon and perovskite cells\autocite{chenOrganicSolarCells2019}.
Much of the advancement is credited to the development of PM6:Y6 donor:acceptor systems\autocite{huEfficientOrganicSolar2026}.
Current challenges in the development of commercially viable materials for OSCs are stability\autocite{tegegneProgressStableOrganic2025} and solubility in green solvents for deposition.
This proposal is the implementation of the Knoevenagel condensation\autocite{ohSynthesisConjugatedPolymer2021} in the preparation of PM6 and Y6 inspired green material libraries and the screening of these libraries for preferred properties of irradiation stability and solubility in green solvents, followed by implementation in bulk heterojunction OSCs to determine PCE.

The Knoevenagel condensation is an aldol condensation that typically forms conjugated systems.
This reaction often has greatest yields in aqueous sodium hydroxide solutions with small quantities of organic solvent present5.
This condensation reaction has a much lower environmental impact than the Stille coupling commonly used to synthesize conjugated materials for OSCs.
A combinatorial library of polymeric donor materials produced via the Knoevenagel polymerization can be screened by irradiation with light, and extraction with green solvents such as 2-methyl tetrahydrofuran.
The surviving polymers can be examined for PCE with multiple NFAs.

The Alex Lippert lab has experience with the Knoevenagel condensation reaction and experience in aynthesis.
This expertise can be applied to production of novelmaterials, and application in spin coating ITO plates for screening candidate compounds and substrates.

\textbf{Aim1} Screening polymeric donor combinatorial libraries for stability by irradiation degradation and green solvent extraction.

\textbf{Aim2} Tin-free synthesis of small molecule and Y6-like NFAs and polymeric PM6-like
donor combinatorial libraries by Knoevenagel condensation in water.

\textbf{Aim3} Implementation of optimized materials into bulk heterojunction OSCs to determine PCE.

\textbf{Intelectual merit} This research will shed light onto novel and green routes to promissing OCSs.
Not only will new insights into screening methodologies be discovered, but the details of how products formed in water can be stable in long term solar cell aplications will also be elucidated.
New applications of previously underutilized knovenagel conductive polymers will be applied and studied for unique and advantagous properties.

\textbf{Broader impact} This research may provide new tools to progess green solar panel production, and continue to open doors to provide a feasible product that can be cheaply mass produced via the roll-to-toll process.
Solar is a promising provider of green energy, necessary for powering the future, but it faces multiple practical challenges, which this research may provide solutions.
